\section{Introduction and Problem Statement}

\paragraph{Problem Statement:} The task of this project is to analyze peptide abundances under different drug concentrations to determine whether a particular peptide is affected by a particular drug. Another goal is to identify groups of peptides that respond to that drug at a given concentration. Mathematically, given a matrix of protein abundance measurements collected across multiple drug concentrations and matched control conditions, we aim at determining for each drug protein pair whether the observed abundance change at one or more concentration levels is statistically inconsistent with experimental variability. On top of this, we can formulate the problem as a machine learning (ML) task by predicting missing drug concentration values for each peptide-drug pair. Motivation is to help in reducing the experimental workload of drug testing.

\paragraph{Input:}   Let $X \in \mathbb{R}^{P\times C}$ be the abundance matrix for a given drug, where P is the number of detected proteins and $C=8$ is the number of concentration levels. Let $B \in \mathbb{R}^{P\times N}$ be the corresponding background matrix derived from the $N=284$ control runs. For our ML task, input is the concatenation of a drug one-hot (50-dim for 50 drugs) and $\log_{10}(\text{concentration} + 1)$ (plus 1 dimension is for the control condition) for each drug–concentration condition in the dataset~\cite{massive081932}.

For handling the missing values under low mass spectrometry values, we will distinguish three cases: (1) whether a peptide is absent in the drug concentration runs, (2) the control run, or (3) both. Peptides absent in concentration runs but present in the control, or vice versa, could carry valuable information while peptides absent in both are uninformative and excluded. For the ML models, peptides detected in more than $80\%$ of conditions across all drugs will be retained in the input vector. Peptides below the $80\%$ threshold are excluded from the input vector entirely.

\paragraph{Output:} The primary statistical inference output is a score or binary hit indicator per drug, peptide, concentration triple. With hypothesis testing methods, we will judge the relation with p-values. Once peptide-level scores are obtained, the key analysis step is aggregating them into group-level perturbation profiles per drug and discovering  how sets of peptides belonging to the same protein, kinase family, or biological pathway are jointly affected. For our ML approach, output is the vector of log2-transformed intensities for a chosen set of about 3,000 well-measured peptides (those present in at least 80\% of conditions across all drugs).


\section{Proposed approach}

\subsection{Simple Statistical Methods}
A naive initial approach is to apply simple statistical methods that model the relationship between drug type, drug concentration, and peptide abundance. More specifically, we can treat each peptide as an output variable and test whether its abundance changes is associated with drug type and concentration. For example, for a specific peptide, we can fit a linear model such as:
\\
\\peptide abundance = drug label + concentration + drug-by-concentration interaction
\\
\\The coefficients of this linear model allow us to assess whether peptide abundance changes based on the concentration of a drug, and whether this relationship differs across drug types. If the concentration–response relationship is not well approximated by a linear model, we can instead fit a more flexible dose–response curve. We then repeat this analysis across all peptides and apply multiple-testing correction, such as false discovery rate control, to account for the large number of simultaneous tests. This gives a basic first-pass map of which peptides appear to respond to which drugs.

\subsection{Multi-task Neural Regressor}
We also propose a multi-task neural regressor that maps each drug–concentration condition to the joint peptide-intensity vector it is expected to produce. Each training example concatenates a 50-dimensional one-hot over the 50 drugs with a one-dimensional dose feature, the base-10 logarithm of (concentration + 1), with DMSO encoded as zero. The target is a roughly 3,000-dimensional vector of log2 intensities for peptides that are present in at least 80 percent of conditions across all drugs, and missing entries are masked out of the loss. We plan to train a small feed-forward network with masked MSE in PyTorch, following the starter-code template. Predicting all peptide intensities jointly from a single (drug, dose) input forces the network to learn cross-peptide correlations, so peptides that vary together under the same drug share statistical strength. This directly answers the handout's challenge to detect joint associations of multiple peptides to a condition.
Because the network is a regressor rather than a target caller, we convert its predictions into a per-drug target ranking using integrated gradients of each predicted peptide intensity with respect to the drug input, evaluated along the trajectory from DMSO to the highest concentration. Summing the absolute attribution across peptides assigned to the same top canonical protein yields a per-protein score for each drug and produces a ranked candidate-target list that we can compare against the Kd table from \cite{klaeger2017target}. Since each drug–concentration condition is observed only once, overfitting is the dominant risk. We plan on addressing this by keeping the network small, applying dropout and early stopping on a drug-wise validation split, and reporting only targets whose attribution direction agrees with the DMSO-versus-maximum-dose log fold change. All hyperparameters are tuned on the validation drugs and never on the test drugs. Peptide variants are prefiltered before training, and we will report how many variants remain after each filter in the preliminary statistics.


\section{Evaluation of the project models}

\subsection{Simple Statistical Methods}

For the simple statistical methods, we will look for target peptides whose abundance shows a consistent dose-response trend under a drug, and we will measure how many such associations remain significant after correcting for multiple comparisons. As a knowledge prior, we can also check whether the strongest discovered associations match prior biological knowledge, such as known targets or pathway members related to that drug, when such information is available.


\subsection{Neural Methods}

For the neural regressor, evaluation has two concurrent goals. First, we plan on establishing that the network generalizes beyond the conditions it was trained on, and establishing that the targets it proposes are biologically meaningful rather than artifacts of noise or overfitting on single replicates. We design three data splits to separate these concerns, and we pair them with metrics that stress both ranking quality and statistical calibration.
To probe generalization, we hold out five of the fifty drugs entirely, train the network on the remaining forty-five, and ask whether its target rankings for the held-out drugs recover their published Kd profiles. This drug-holdout set is our primary generalization test. Within the training drugs, we additionally withhold one concentration per drug (for example, 300 nM) as a dose-holdout validation set. This set is used for hyperparameter tuning and early stopping and measures whether the regressor generalizes across doses it has not seen. Finally, we pool all fifty DMSO controls to construct a null distribution of peptide intensities under no drug. This DMSO-only background could be used to calibrate significance thresholds both for the NN's attribution scores and for the Welch t-test baseline introduced in Section 2.1.
Our primary metric is the recovery of known kinase targets. The table from \cite{klaeger2017target} provides dissociation constants between each drug and its known kinase binders, and we treat kinases with published Kd below one micromolar as true positives per drug and all other proteins as negatives. Proteins are then ranked by the network's per-drug attribution score, and we report precision at k for k equal to 1, 5, and 10, together with the area under the precision-recall curve (AUPRC), macro-averaged across the held-out drugs. Precision at k captures whether the top of the ranking is dominated by biologically correct targets, which is what matters in practice because a drug usually has only a handful of relevant kinase targets. AUPRC captures ranking quality across the full list and could be more robust than AUROC under the severe class imbalance expected here, where true target kinases form on the order of one percent of all detected proteins per drug.
We add to the primary recovery metric with three diagnostic checks. First, we would like to calibrate the pipeline on permuted DMSO data by randomly relabeling DMSO controls as treatment versus control, running the full attribution pipeline, and verifying that the resulting p-value distribution is approximately uniform.  we report the empirical versus nominal false positive rate to check for any systematic bias. Second, we assess dose response plausibility by fitting a Hill curve through the treated intensities of each called target and retaining only hits whose curve is monotonic and whose fitted EC50 falls inside the sampled concentration window of 3 nM to 30 µM, discarding non-monotonic hits as noise. Third, we would like to check cross-drug consistency by computing the Jaccard overlap of top-10 target calls for drug pairs that share at least one known kinase target according to \cite{klaeger2017target}, reasoning that drugs with overlapping published kinase profiles should produce overlapping rankings if the method is picking up real pharmacology rather than memorizing drug identity.



\section{Assessing limitations and alternative approaches}

The proposed approach for identifying the relationship between protein type and drug concentration may be affected by limitations in the input data. One example is the curse of dimensionality, where the relatively small number of samples for each protein/drug concentration combination compared to the number of peptide variants measured. When the number of features greatly exceeds the number of samples, overfitting can happen in predictive models and make it more likely to identify false relationships between drug treatment and peptide abundance, especially in a neural regressor model. Another limitation is the potential of the variants intensity file having missing values or zero entries. Because these values might not occur randomly, due to detection limits or variability in mass spectrometry, they might introduce bias and lead to uncertain estimates of peptide abundance changes caused by drug treatment. A final important limitation of simple statistical analysis on peptides is that it may assume all peptide changes from drug treatment are independent. In reality, a drug might affect a certain peptide which, in turn, affects other peptides. This type of indirect relationship between peptide abundances would not be captured by a model considering each peptide as independent. Similarly, peptides from the same proteins may show similar responses to drugs, which further shows that independence assumptions may not be valid in tests such as the Student t-test.

A few strategies can be used to reduce the impact of these limitations. The first strategy is aggregating peptides to their protein-level summaries in order to reduce redundancy and strengthen the understanding of changes done by the drug. Regularized regression methods can also be used to mitigate overfitting caused by missing values. A mitigation for false discoveries is multiple testing correction procedures when we perform large-scale abundance testing. 
% \Jackie{ [add sentence about neural regressor limitation]}

